\section{Count Subsets with Sum K}
\label{sec:dp-count-subsets-sum-k}

\problemstatement{Given an array $\textit{arr}$ of non-negative integers
and a target $k$, count \textbf{how many} subsets sum to exactly $k$.}

\begin{patternbox}
This is Section~\ref{sec:dp-subset-sum-target}'s exact recurrence with
the logical OR replaced by \textbf{addition} -- the same
exists-to-count transformation as the Recursion chapter's
``Check if There Exists a Subsequence with Sum K'' to ``Count
Subsequences with Sum K'' transition,
now applied on top of the memoized/tabulated table instead of plain
recursion.
\end{patternbox}

\subsection*{Deriving the recurrence}

\[
  f(\textit{index}, \textit{target}) = f(\textit{index}{-}1, \textit{target})
  \ +\ \big(\textit{arr}[\textit{index}] \le \textit{target}\ ?\
  f(\textit{index}{-}1, \textit{target}-\textit{arr}[\textit{index}]) : 0\big).
\]
Base case (handling a possible zero-valued element carefully, since a
$0$ can be included or excluded without changing the sum, doubling the
count): if $\textit{index}=0$, return $2$ if $\textit{target}=0$ and
$\textit{arr}[0]=0$; return $1$ if $\textit{target}=0$ or
$\textit{target}=\textit{arr}[0]$ (but not both conditions above); else
return $0$.

\subsection*{Approach 1: Recursion}

\begin{lstlisting}
#include <bits/stdc++.h>
using namespace std;

int countSubsetsRecursive(int index, int target, vector<int>& arr) {
    if (index == 0) {
        if (target == 0 && arr[0] == 0) return 2;
        if (target == 0 || target == arr[0]) return 1;
        return 0;
    }

    int notPick = countSubsetsRecursive(index - 1, target, arr);
    int pick = 0;
    if (arr[index] <= target) {
        pick = countSubsetsRecursive(index - 1, target - arr[index], arr);
    }
    return notPick + pick;
}

int main() {
    vector<int> arr = {1, 2, 2, 3};
    int n = arr.size();
    cout << countSubsetsRecursive(n - 1, 3, arr) << endl; // 3
    return 0;
}
\end{lstlisting}

\begin{complexitybox}[Approach 1 Complexity]
\textbf{Time.} $O(2^n)$.
\textbf{Space.} $O(n)$ recursion stack.
\end{complexitybox}

\subsection*{Approach 2: Memoization}

\begin{lstlisting}
#include <bits/stdc++.h>
using namespace std;

int countSubsetsMemo(int index, int target, vector<int>& arr,
                      vector<vector<int>>& memo) {
    if (index == 0) {
        if (target == 0 && arr[0] == 0) return 2;
        if (target == 0 || target == arr[0]) return 1;
        return 0;
    }
    if (memo[index][target] != -1) return memo[index][target];

    int notPick = countSubsetsMemo(index - 1, target, arr, memo);
    int pick = 0;
    if (arr[index] <= target) {
        pick = countSubsetsMemo(index - 1, target - arr[index], arr, memo);
    }
    return memo[index][target] = notPick + pick;
}

int main() {
    vector<int> arr = {1, 2, 2, 3};
    int n = arr.size(), k = 3;
    vector<vector<int>> memo(n, vector<int>(k + 1, -1));
    cout << countSubsetsMemo(n - 1, k, arr, memo) << endl; // 3
    return 0;
}
\end{lstlisting}

\begin{complexitybox}[Approach 2 Complexity]
\textbf{Time.} $O(n \cdot k)$.
\textbf{Space.} $O(n \cdot k)$ memo $+$ $O(n)$ recursion stack.
\end{complexitybox}

\subsection*{Approach 3: Tabulation}

\begin{lstlisting}
#include <bits/stdc++.h>
using namespace std;

int countSubsetsTabulation(vector<int>& arr, int k) {
    int n = arr.size();
    vector<vector<int>> dp(n, vector<int>(k + 1, 0));

    dp[0][0] = (arr[0] == 0) ? 2 : 1;
    if (arr[0] != 0 && arr[0] <= k) dp[0][arr[0]] = 1;

    for (int index = 1; index < n; index++) {
        for (int target = 0; target <= k; target++) {
            int notPick = dp[index - 1][target];
            int pick = (arr[index] <= target) ? dp[index - 1][target - arr[index]] : 0;
            dp[index][target] = notPick + pick;
        }
    }
    return dp[n - 1][k];
}

int main() {
    vector<int> arr = {1, 2, 2, 3};
    cout << countSubsetsTabulation(arr, 3) << endl; // 3
    return 0;
}
\end{lstlisting}

\subsection*{Dry run of the tabulation table}

\dryrun{Take $\textit{arr} = [1,2,2,3]$, $k=3$.}

\begin{center}
\begin{tabular}{c|c|c|c|c}
$\textit{index}\backslash\textit{target}$ & 0 & 1 & 2 & 3 \\ \hline
$0$ (val $1$) & 1 & 1 & 0 & 0 \\
$1$ (val $2$) & 1 & 1 & 1 & 1 \\
$2$ (val $2$) & 1 & 1 & 2 & 2 \\
$3$ (val $3$) & 1 & 1 & 2 & 3
\end{tabular}
\end{center}

The three qualifying subsets summing to $3$ are $\{1,2\}$ (using the
first $2$), $\{1,2\}$ (using the second $2$), and $\{3\}$ alone --
matching $\textit{dp}[3][3]=3$.

\begin{complexitybox}[Approach 3 Complexity]
\textbf{Time.} $O(n \cdot k)$.
\textbf{Space.} $O(n \cdot k)$ table, no recursion stack.
\end{complexitybox}

\subsection*{Approach 4: Space Optimization}

\begin{lstlisting}
#include <bits/stdc++.h>
using namespace std;

int countSubsetsSpaceOptimized(vector<int>& arr, int k) {
    int n = arr.size();
    vector<int> prev(k + 1, 0);

    prev[0] = (arr[0] == 0) ? 2 : 1;
    if (arr[0] != 0 && arr[0] <= k) prev[arr[0]] = 1;

    for (int index = 1; index < n; index++) {
        vector<int> current(k + 1, 0);
        for (int target = 0; target <= k; target++) {
            int notPick = prev[target];
            int pick = (arr[index] <= target) ? prev[target - arr[index]] : 0;
            current[target] = notPick + pick;
        }
        prev = current;
    }
    return prev[k];
}

int main() {
    vector<int> arr = {1, 2, 2, 3};
    cout << countSubsetsSpaceOptimized(arr, 3) << endl; // 3
    return 0;
}
\end{lstlisting}

\begin{complexitybox}[Approach 4 Complexity]
\textbf{Time.} $O(n \cdot k)$.
\textbf{Space.} $O(k)$ for the rolling row.
\end{complexitybox}

\begin{takeawaybox}
Counting subsets and checking subset existence differ by exactly one
operator (addition instead of logical OR) at every step of an otherwise
identical recurrence -- the same lesson already drawn from the
Recursion chapter's exists-to-count transition.
Watch for the zero-value edge case specifically whenever a problem allows
zero-valued array elements, since a $0$ can always be freely included or
excluded without affecting the sum, silently doubling some counts if not
handled in the base case.
\end{takeawaybox}
